\documentclass[a4paper, 11pt]{article}
\usepackage[top=2.1cm, bottom=2.1cm, left=2.3cm, right=2.3cm]{geometry}
\usepackage[utf8]{inputenc}
\usepackage[T1]{fontenc}
\usepackage[italian]{babel}
\usepackage{amsmath, amssymb}
\usepackage{siunitx}
\usepackage{setspace}
\usepackage{titlesec}
\usepackage[colorlinks=true, urlcolor=blue, linkcolor=black]{hyperref}
\usepackage{lmodern}
\usepackage{microtype}

\setstretch{1.14}
\setlength{\parindent}{15pt}
\setlength{\parskip}{0.5em}


\titleformat{\section}
  {\normalfont\large\bfseries}
  {\thesection}{1em}{}
\titlespacing*{\section}{0pt}{1.5ex plus .5ex minus .2ex}{0.8ex plus .1ex}

\pagestyle{plain}

\begin{document}

\vspace{-0pt}
\noindent \textit{\small{Dipartimento di Fisica ``G. Occhialini'' - Corso di Laurea triennale in Fisica}}\\
\vspace{-18pt}
% \small{Organización y/o Universidad}}
\begin{center}
\normalsize{Riassunto dell'elaborato:}\\
\medskip
\Large{``Metodo dei fasori applicato all'analisi dei tempi di vita di fluorescenza''}
\vspace{0.3cm}
\normalsize

\noindent
\begin{tabular*}{\textwidth}{@{\extracolsep{\fill}}ll@{}}
    \textbf{Candidata:} Sara Di Lernia & \textbf{Relatrice:} Prof.ssa Laura Sironi \\
    \textbf{Matricola:} 898437 & \textbf{Correlatore:} Dott. Davide Panzeri \\
    \textbf{Recapito telefonico:} +39 345 830 7326 & \textbf{Appello di Laurea:} Novembre 2026 \\
\end{tabular*}

\vspace{0.2cm}
\hrule
\vspace{0.4cm}


\normalsize

\end{center}

\vspace{-5pt}

\noindent L'obiettivo principale di questo studio è quello di approfondire il \textit{metodo dei fasori} per l'analisi dei dati FLIM (Fluorescence Lifetime Imaging Microscopy), prima mediante la generazione e l'analisi di decadimenti di fluorescenza simulati, poi tramite lo studio di acquisizioni reali. Infine viene presentata un'interfaccia grafica interattiva basata sul codice implementato.

La microscopia a tempi di vita di fluorescenza (FLIM) è una tecnica di imaging spettroscopica che permette di mappare i tempi di vita delle specie fluorescenti presenti in un campione biologico. Il contrasto delle immagini ottenute, quindi, non dipende dall'intensità della fluorescenza ma dal tempo medio $\tau$ che una molecola trascorre nello stato eccitato prima di tornare allo stato fondamentale. Questo decadimento spontaneo avviene secondo una legge esponenziale $I(t)=I_0\,e^{-t/\tau}$.\\
Il tempo di vita di fluorescenza è sensibile alle variazioni ambientali, come pH e stato ossidativo, dunque la FLIM può essere utilizzata anche per monitorare processi biologici come il trasferimento di energia per risonanza di fluorescenza (FRET), o la concentrazione relativa di NADH e FAD, due coenzimi coinvolti nel metabolismo cellulare.\\
L'acquisizione sperimentale di immagini FLIM avviene tipicamente mediante la tecnica TCSPC (Time-Correlated Single Photon Counting): il campione viene scansionato pixel per pixel sfruttando un laser a impulsi periodici e molto brevi ($\sim10^{-12}\,\text{s}$, $f\sim 20-80\,\text{MHz}$), mentre un rivelatore registra i fotoni di fluorescenza emessi e il loro tempo di ritardo rispetto all'impulso di eccitazione.
Ripetendo la misura per migliaia di impulsi, per ogni pixel viene accumulato un istogramma dei ritardi, che ricostruisce la curva di decadimento di fluorescenza.

L'analisi FLIM tradizionale si basa sull'interpolazione esponenziale o multi-esponenziale del decadimento di fluorescenza raccolto da ciascun pixel dell'immagine, risultando computazionalmente costosa e spesso imprecisa a causa del numero di parametri liberi.\\
Un approccio più moderno è invece quello dell'analisi dei fasori: ogni decadimento viene trasposto in un punto del piano complesso $z(\omega) = G(\omega) + i\,S(\omega)$ mediante la trasformata di Fourier del segnale di fluorescenza calcolata alla frequenza angolare di modulazione ($\omega =2\pi/T_{laser}$), e normalizzata sull'intensità totale emessa nel pixel:
\vspace{-5pt}
\[
G(\omega) = \frac{1}{I_{tot}} \int_0^{T} I(t) \cos(\omega t)\,dt, \qquad 
S(\omega) = \frac{1}{I_{tot}} \int_0^{T} I(t) \sin(\omega t)\,dt, \qquad 
\text{con }\;I_{tot} = \int_0^{T} I(t)\,dt
\]
\vspace{-22pt}

\noindent In questo spazio, i punti associati ai decadimenti mono-esponenziali puri e senza rumore cadono su una curva di riferimento detta ``semicerchio universale'', mentre quelli multi-esponenziali (emissione di specie fluorofore con tempi di vita diversi) si trovano nell'area compresa tra questa curva e l'asse reale. Grazie alla linearità della trasformata di Fourier, i fasori obbediscono alla regola della somma vettoriale pesata sull'intensità frazionaria dei singoli fluorofori; questo permette di effettuare la scomposizione multicomponente secondo regole geometriche.\\
Nei dati FLIM reali, i decadimenti di fluorescenza sono convoluti temporalmente con la funzione di risposta dello strumento (IRF, \textit{Instrument Response Function}), fortemente sensibile alla strumentazione utilizzata. Nell'analisi dei fasori, è possibile sfruttare il teorema della convoluzione per isolare il segnale di fluorescenza dalla IRF mediante una semplice divisione complessa sul piano dei fasori: $z_{\text{decadimento}} = z_{\text{segnale}}/z_{\text{IRF}}$.\\
La deconvoluzione del segnale dalla IRF tramite l'approccio appena presentato è stata studiata sia su dati simulati che reali.

In una prima fase sono state generate delle simulazioni in linguaggio Python (sfruttando principalmente le librerie \textit{Phasorpy}, \textit{Numpy}, \textit{Matplotlib}), volte allo studio delle funzioni di analisi delle curve di decadimento.
In primo luogo sono stati simulati decadimenti mono-esponenziali e multi-esponenziali (a due e tre componenti) al variare di $\tau_k$ e $f_k$ (tempi di vita e intensità frazionarie di ogni componente), in seguito rappresentati sul piano dei fasori. Al crescere di $\tau$ i fasori si spostano dall'estremo destro del semicerchio in $(1, 0)$ verso quello sinistro in $(0, 0)$, quest'ultimo corrispondente al valore massimo di $12.5\, \text{ns}$, pari al periodo del laser. È stato possibile verificare come il fasore di un decadimento dato da una combinazione di due specie giaccia sul segmento congiungente le due specie pure, a una distanza da ciascun estremo che è inversamente proporzionale alla $f_k$ della specie corrispondente.
Analogamente, per i decadimenti a tripla esponenziale, è stata verificata la posizione del fasore nell'area del triangolo che ha per vertici le specie pure; la $f_k$ di ciascuna componente risulta proporzionale all'area del sotto-triangolo opposto al vertice corrispondente. Infine è stato studiato come l'aggiunta di rumore gaussiano sposti sistematicamente il fasore relativo ai decadimenti verso il centro del semicerchio.

Nella sezione principale della trattazione, gli algoritmi di analisi sviluppati sono stati applicati a due set di immagini FLIM reali, estratte dall'archivio open-source Zenodo (progetto FLUTE).\\
In entrambe le immagini, il fasore della IRF è stato estratto grazie a immagini di calibrazione con fluorofori a $\tau$ noto, come la fluoresceina ($\tau=4\,\text{ns}$).\\
Per quanto riguarda il primo campione biologico analizzato, sono stati rappresentati i fasori associati ai singoli pixel, dopo aver applicato un filtro \textit{median} \texttt{3}$\times$\texttt{3} alla matrice delle coordinate $(G, S)$, per ridurre la dispersione dei fasori causata dal rumore. È inoltre stata introdotta una maschera per la selezione spaziale di sezioni dell'immagine (ROI, Region Of Interest), che permette di proiettare sul piano dei fasori solo i pixel di interesse.\\
La seconda immagine invece era corredata da due acquisizioni ottenute filtrando rispettivamente l'emissione nel rosso ($\lambda = 607\, \text{nm}$) e nel verde ($\lambda = 550\, \text{nm}$) di due specie fluorofore differenti. Questo ha permesso di ottenere dei riferimenti per effettuare la scomposizione delle frazioni di intensità sul piano dei fasori.

Selezionate due ROI distinte sull'immagine, è stato determinato il rapporto di intensità dei fluorofori rosso e verde in ciascuna regione, valutando la posizione relativa dei fasori medi di ogni ROI rispetto ai riferimenti. Le incertezze relative alla varianza dei cluster sono state propagate ai risultati mediante simulazione Monte Carlo.
I risultati ottenuti mostrano, nella prima regione, una predominanza della componente rossa ($f_{\text{rosso}} = 0.94 \pm 0.04$ e $f_{\text{verde}} = 0.06 \pm 0.04$), mentre nella seconda si osserva una presenza bilanciata di entrambe le componenti ($f_{\text{rosso}} = 0.55 \pm 0.12$ e $f_{\text{verde}} = 0.45 \pm 0.12$).


In conclusione, da questa trattazione emerge come l'analisi dei fasori rappresenti una valida alternativa al \textit{fitting} non lineare, riducendo i costi computazionali legati all'interpolazione e semplificando il processo di deconvoluzione dalla IRF.\\
L'interfaccia grafica (sviluppata mediante la libreria \textit{Panel}) ha lo scopo di facilitare la visualizzazione sia delle fasi di generazione dei segnali, sia della corrispondenza spaziale tra le distribuzioni sul piano dei fasori e i pixel di un'immagine FLIM, agevolando l'introduzione all'analisi dei fasori ai non addetti al settore.
\end{document}
